\chapter{Различные операторы.}
\section{Сопряжённые операторы.}\label{seq:conjugate-operators}
\begin{definition}
    Пусть $X, Y$ -- нормированные пространства и $A \in \LL(X, Y)$.
    Оператор $A^*\colon Y^* \to X^*$, определённый как
    \[
        (A^*g)(x) = g(Ax),
    \]
    будем называть сопряжённым к $A$.
\end{definition}
\begin{note}
    Заметим, что $g(Ax) = (g \circ A)(x)$ при $g \in Y^*$ и $x \in X$.
    То есть
    \[
        \forall g \in Y^*\colon A^*g = g \circ A.
    \]
    Поскольку $g \in \LL(Y, \Cm)$ и $A \in \LL(X, Y)$, то $g \circ A \in \LL(X, \Cm) = X^*$, поэтому определение корректно.

    Очевидно, что $A^*$ -- линеен:
    \[
        A^*(\alpha_1 g_1 + \alpha_2 g_2)(x) = (\alpha_1 g_1 + \alpha_2 g_2)(Ax) = \alpha_1 g_1(Ax) + \alpha_2 g_2 (Ax) = \alpha_1 (A^*g_1)(x) + \alpha_2 (A^*g_2)(x).
    \]
    Следовательно, $A^*(\alpha_1 g_1 + \alpha_2 g_2) = \alpha_1 A^* g_1 + \alpha_2 A^* g_2$.

    Найдём норму $A^*$.
    \[
        \|A^*\| = \sup_{\|g\| \leq 1} \|A^*g\| = \sup_{\|g\| \leq 1} \sup_{\|x\| \leq 1} |g(Ax)| = \sup_{\|x\| \leq 1}\sup_{\|g\| \leq 1} |g(Ax)| = \sup_{\|x\| \leq 1} \|Ax\| = \|A\|,
    \]
    где предпоследний переход сделан в силу следствия из теоремы Хана-Банаха.

    Таким образом $\|A^*\| = \|A\| < +\infty$ и $A^* \in \LL(Y^*, X^*)$.
\end{note}
\begin{note}
    Рассмотрим второй сопряжённый оператор $A^{**}\colon X^{**} \to Y^{**}$.
    Пусть $\Phi_X\colon X \to X^{**}$ и $\Phi_Y\colon Y \to Y^{**}$ -- канонические изометрии Банаха.

    Заметим, что $\forall x \in X$ и $g \in Y^*\colon$
    \[
        (A^{**}\Phi_X(x))(g) = \Phi_X(x)(A^* g) = (A^* g)(x) = g(Ax) = (\Phi_Y(Ax))(g).
    \]
    А значит $A^{**}\Phi_X(X) \subset \Phi_Y(Y)$.
    Поскольку на $\Phi_Y(Y)$ каноническая изометрия $\Phi_Y$ -- обратима, то
    \[
        \Phi_Y^{-1} \circ A^{**} \circ \Phi_X = A.
    \]
\end{note}
\begin{proposition}
    Если $A \in \LL(X, Y)$ и $\exists A^{-1}\in \LL(Y, X)$, то $\exists (A^*)^{-1} \in \LL(X^*, Y^*)$ и, более того, $(A^*)^{-1} = (A^{-1})^*$.
\end{proposition}
\begin{proof}
    Рассмотрим $T = (A^{-1})^* \in \LL(X^*, Y^*)$.
    Запишем действие $A^*$ на $T\colon$
    \[
        (A^*(Tf))(x) = (Tf)(Ax) = ((A^{-1})^* f)(Ax) = f(A^{-1}A x) = f(x).
    \]
    А значит
    \[
        A^* T = I,
    \]
    то есть $T$ -- правый обратный к $A^*$.
    Аналогично
    \[
        T(A^* g)(y) = ((A^{-1})^* A^* g)(y) = (A^* g)(A^{-1}y) = g(AA^{-1}y) = g(y),
    \]
    и $T$ является и левым обратным, то есть $TA^* = I$.
    Значит $T$ -- обратный к $A^*$ и, следовательно, $(A^*)^{-1} = (A^{-1})^*$.
\end{proof}
\begin{note}
    Пусть $A \in \LL(X, Y)$ и существует $(A^*)^{-1} \in \LL(X^*, Y^*)$.
    Тогда, по предыдущему утверждению, существует $(A^{**})^{-1} = ((A^*)^{-1})^* \in \LL(Y^{**}, X^{**})$.
    С другой стороны, так как
    \[
        A = \Phi_Y^{-1} \circ A^{**} \circ \Phi_X,
    \]
    и если $x \in \ker A$, то $\Phi_X(x) \in \ker A^{**}$, то $\ker A = 0$.
    Тогда на образе у $A$ существует обратный $A^{-1}\colon \im A \to X$, и он имеет вид
    \[
        A^{-1} = \Phi_X^{-1} \circ (A^{**})^{-1} \circ \Phi_Y.
    \]
    Более того, $A^{-1} \in \LL(\im A, X)$, поскольку (в силу изометричности $\Phi_X$ и $\Phi_Y$)
    \[
        \|A^{-1} y\| \leq \|(A^{**})^{-1}\| \|y\|
    \]
\end{note}
\section{Теорема Фредгольма.}
\begin{theorem}[Фредгольм]\label{th:fredhgolm}
    Пусть $A \in \LL(X, Y)$.
    Тогда верно следующее:
    \[
        \begin{cases}
            \ker A =  ^{\perp}\im A^*, \\
            \ker A^* = \im A^{\perp}.
        \end{cases}
    \]
    И справедливы такие соотношения:
    \[
        \begin{cases}
        (^\perp(\im A^*))^{\perp} = [\im A^*]_{*w}, \\
        ^\perp((\im A)^\perp) = [\im A]_{\|\cdot\|}.
        \end{cases}
    \]
\end{theorem}
\begin{proof}
    Пусть $g \in \ker A^* \subset Y^*$, то есть $A^* g = 0$.
    А значит
    \[
        (A^* g)(x) = 0, \forall x \in X.
    \]
    По определению сопряжённого оператора это эквивалентно
    \[
        g(Ax) = 0, \forall x \in X.
    \]
    то есть, по определению, $g \in \im A^\perp$.

    Пусть $x \in \ker A \subset X$, то есть $Ax = 0$.
    Это означает, что $\forall f \in Y^* f(Ax) = 0$.
    Или, в силу определения сопряжённого оператора
    \[
        (A^* f)x = 0.
    \]
    Но, по определению аннулятора, это и означает $x \in ^\perp \im A^*$.

    Вторые два равенства получаются применением следствий для аннуляторов к $\im A$ и $\im A^*$.
\end{proof}
\begin{theorem}\label{th:closability}
    Пусть $X, Y$ -- банаховы пространства и $A \in \LL(X, Y)$.
    Пусть $\im A$ -- сильно замкнут.
    Тогда $\im A^* = \ker A^\perp$ и $\im A = {}^\perp \ker A^*$.
\end{theorem}
\begin{proof}
    Второе равенство тривиально следует из замкнутости образа $A$ и ранее полученных соотношений для аннуляторов.

    Достаточно показать, что $\ker A ^\perp \subset \im A^*$, поскольку в силу теоремы Фредгольма
    \[
        \im A^* \subset [\im A^*]_{*w} = \ker A^{\perp}.
    \]
    Это означает, что $\forall f \in X^*\colon \forall x \in \ker A \hookrightarrow f(x) = 0$ существует $g \in Y^*$ такой, что $f = A^* g$.
    По определению сопряжённого оператора
    \[
        f(x) = (A^* g)(x) = g(Ax).
    \]
    Введём функционал $\phi\colon \im A \to \Cm$ как
    \[
        \phi(Ax) = f(x).
    \]
    Определение корректно, поскольку если $Ax_1 = Ax_2$, то $x_1 - x_2 \in \ker A$, а значит $f(x_1 - x_2) = 0$ -- из чего и следует $f(x_1) = f(x_2) = \phi(Ax_1) = \phi(Ax_2)$.

    Поскольку $\im A$ -- замкнутое подпространства банахового $Y$, то оно само по себе является банаховым пространством.
    Поэтому $A\colon X \to \im A$ -- открытое отображение и, следовательно, существует $r > 0$:
    \[
        O_r^Y \cap \im A \subset A(O_1^X(0)).
    \]
    Покажем непрерывность $\phi$.
    Поскольку $\forall y \in \im A \setminus \{0\}$ верно, что
    \[
        \frac{r}{2}\frac{y}{\|y\|} \in O_r^Y(0) \cap \im A \subset A(O_1^X(0)),
    \]
    то существует $x \in O_1^X(0)$, образом которого является $\frac{r}{2}\frac{y}{\|y\|}$:
    \[
        Ax = \frac{r}{2}\frac{y}{\|y\|}.
    \]
    Следовательно
    \[
        y = A\left(\frac{2\|y\|}{r}x\right),
    \]
    то есть
    \[
        \phi(y) = \phi\left(A\left(\frac{2\|y\|}{r}x\right)\right) = \frac{2\|y\|}{r}f(x).
    \]
    Теперь мы можем оценить норму $\phi$:
    \[
        |\phi(y)| \leq \frac{2\|y\|}{r}\|f\| \Rightarrow \|\phi\| \leq \frac{2\|f\|}{r}.
    \]
    Таким образом $\phi$ -- непрерывен на $\im A$, и мы можем продолжить его на всё $Y$ с сохранением нормы.
    Пусть $g\in Y^*\colon g|_{\im A} = \phi, \ \|g\|=\|\phi\|$.
    Так как
    \[
        f(x) = \phi(Ax) = g(Ax) = A^* g(x),
    \]
    то $f = A^* g$ и $f \in \im A^*$ -- что и завершает доказательство теоремы.
\end{proof}
\begin{example}
    Банаховость $Y$ существенна.

    Рассмотрим оператор $A$, определяемый как
    \[
        A\colon (\ell^1, \|\cdot\|_1) \to (\ell^1, \|\cdot\|_2), \ x \mapsto x.
    \]
    Очевидно, что $\im A = \ell^1$, и, таким образом его образ сильно замкнут в $(\ell^1, \|\cdot\|_2)$.
    С другой стороны
    \[
        A^*\colon (\ell^2, \|\cdot\|_2) \to (\ell^\infty, \|\cdot\|_\infty),\  y \mapsto y.
    \]
    Заметим, что $\im A^* = \ell^2$.
    Но ядро $A$ -- тривиально, а значит $\ker A ^\perp = \ell^\infty$.

\end{example}
